\section{Scale: the winning baseline changes identity}
\label{sec:scale}

The standing objection to everything above is that 0.5B and 1B models are small
enough that RLVR is not expected to help, making the negative result
unsurprising. We test that objection rather than concede it: the same
CodeTAT-QA task, the same protocol, the same two budgets, three seeds, on
Qwen2.5-7B-Instruct.

\begin{table}[t]
\centering
\small
\setlength{\tabcolsep}{3pt}
\caption{Frozen-test accuracy (\%) at 7B, all three seeds. The solver is now
15 points behind and irrelevant; continued SFT attains the maximum in all six
groups and GRPO trails it in all six.}
\label{tab:scale}
\begin{tabular}{rr rrrrr r}
\toprule
$m$ & seed & solver & $W_m$ & c-SFT & RS-SFT & GRPO & $\Delta$ \\
\midrule
2  & 0 & 64.54 & 71.28 & \textbf{79.79} & 75.18 & 76.60 & $-3.19$ \\
2  & 1 & 64.54 & 69.86 & \textbf{77.66} & 76.95 & 74.82 & $-2.84$ \\
2  & 2 & 64.54 & 70.57 & \textbf{76.95} & 74.82 & 75.18 & $-1.77$ \\
64 & 0 & 64.54 & 79.08 & \textbf{81.91} & 80.50 & 79.79 & $-2.13$ \\
64 & 1 & 64.54 & 78.72 & \textbf{80.50} & 80.14 & 79.43 & $-1.06$ \\
64 & 2 & 64.54 & 79.43 & \textbf{85.11} & 81.21 & 80.85 & $-4.26$ \\
\bottomrule
\end{tabular}
\end{table}

Table~\ref{tab:scale} shows the transition. The solver's accuracy is unchanged
by construction---$64.54\%$, the same number as at 0.5B, because it contains no
model---but the learning arms have moved past it. Warm-start SFT on two
demonstrations now reaches $69.86$--$71.28\%$, and every arm at $m{=}64$ is
around $80\%$. The crossing point predicted by the dose-response curves of
Section~\ref{sec:results} is real and lies between 1B and 7B on this task. The
``gains emerge with capability'' account is correct about the solver.

It is not correct about $\Delta$. In all six groups the maximum is attained by
\textbf{continued SFT}---$76.95$--$85.11\%$---and GRPO trails it by $1.06$ to
$4.26$ points (mean $-2.54$). Scaling the backbone $14\times$ replaced the arm
that beats GRPO; it did not produce a setting in which GRPO leads.

\paragraph{What this does to the claim.}
Within the registered budget it generalizes it. Had the solver kept winning at
7B, a reader could reasonably suspect we had over-engineered one baseline for
one regime. Instead the solver exits exactly where the scaling account says it
should, and the arm that replaces it is continued SFT---an arm present in
essentially every RLVR paper, requiring no engineering, and impossible to
characterize as an exotic construction. The methodological requirement that
survives both regimes is not ``include a solver'' but ``make the control group
broad enough that you would notice, and aggregate it with a $\max$.''

\paragraph{But this configuration does not answer the question it was added for,
and we should say so before a reader finds out.}
Config~3 exists to test the objection that 0.5B and 1B are simply too small. It
cannot settle that objection, because at 7B the registered wall-clock budget
buys an order of magnitude less data than it does at 0.5B: $0.13$ epochs over
the prompt pool against $0.53$--$0.79$. We re-ran this section's
smallest-margin cell---$m{=}64$, seed~1, the one most favourable to GRPO---at
ten times the coverage, and it reverses decisively, to $\Delta = +6.74$ and
$+6.38$ across two training seeds, clearing the pre-registered Go threshold
(Appendix~\ref{app:followup}). At the registered budget the result holds,
including under a run that removes the zero-advantage degeneracy we diagnose in
Section~\ref{sec:analysis}. \textbf{What Config~3 demonstrates is therefore a
budget boundary, not a capability one.} We keep it in the paper for that
reason---knowing where a null result stops is worth more than a scale ladder
that does not hold---but the quantitative weight of this study rests on the 11
groups at 0.5B and 1B, not on these six.

\paragraph{Statistical caveats.}
Effect sizes at 7B are smaller than at 0.5B/1B: $-1.06$ to $-4.26$ points on
$n{=}282$ is 3 to 12 items, and the three learning arms lie within a few points
of one another, so which of them attains $\max \mathcal{B}$ is itself sensitive
to sampling noise. Within the registered budget the 7B groups accordingly carry
a qualitative claim---the arm that beats GRPO changes identity, and GRPO is not
that arm in any of the six---rather than an effect size, and past that budget
they do not carry even that. The quantitative result rests on the 11
groups at 0.5B and 1B, where $\max \mathcal{B}$ is attained by a deterministic
solver that leads the nearest learning arm by \emph{at least} $8.5$ points in
Config~1 ($8.51$ to $36.88$ points across the six groups; the quoted figure is
the smallest, not the typical): there the identity of the maximum is not in
question, and no selection over arms enters the estimate.

The largest 7B margin ($-4.26$ points, Config~3 $m{=}64$ seed~2) has an
uncorrected McNemar $p$ of $0.017$. Under the Holm correction this paper
declares across its family of 17 groups, that value becomes $0.254$ and is
\textbf{not} significant; the conference version of this paper quoted the
uncorrected figure alongside a stated Holm procedure, which was inconsistent.
No 7B group is significant after correction. The 7B claim is the qualitative
one stated above, supported by six of six groups agreeing in sign
(Section~\ref{sec:statistics}), and by nothing stronger.
